\begin{lemma}[Typed guarded root forcing and lift equivalence]
\label{lem:typed-guarded-lift-equivalence}
Let \(C\) be as in \cref{con:typed-guarded-lift}, and suppose the
construction reaches \(G_{\mathrm{typ}}(C)\) rather than stopping at a
pairwise profile conflict.  Every visible completion of
\(G_{\mathrm{typ}}(C)\) canonically realises all displayed positive typed
hidden data and all transition-closed root identities.  Moreover,
\(C\) has a hidden completion if and only if
\(G_{\mathrm{typ}}(C)\) has a visible completion (with the distinguished
old higher-support factors kept fixed in the support-normal finite-arity
case).  In particular, hidden non-completion implies visible guarded-lift
non-completion.
\end{lemma}

\begin{proof}
Consider one constrained witness pair with exact domain
\(D=\bigcap_iP_i\), and let \(c\) be its completed atomic pair type.  Every
branch of its private theta bundle remains a specified path, so the typed
path-profile observation gives \(c\in P_i\) for all \(i\), and hence
\[
 c\in D. \tag{TGR}\label{eq:typed-guarded-root-colour}
\]
Thus \(c\) lies in every positive profile imposed on this pair.  Membership
in \(J\) forces a displayed projection equality, and membership in the target
equivalence forces an upward quotient equality.  By
\cref{lem:typed-cylinder-realisation}, membership in a named mixed cylinder
forces precisely the displayed semantic-cylinder relation.
By \cref{lem:typed-transition-cylinder}, membership in the
IDENTITY-FREE positive profile
\(J^\circ\starprod R\starprod(J^R)^\circ\)
forces the NONTRIVIAL named transition
equation for \(\pi_{R,J}\).
Any identity-containing \(R\)
has already reduced to an existing
unary identity term by
\cref{lem:typed-identity-transport-redundant}
and does NOT create a separate
strong theta constraint.
A singleton profile preserves a specified old real pair
type.

It remains to check the additional root-signature equations.  If
\(DD^{\smile}\subseteq J\), then the atom \(c\in D\), the exact
relation \(D\), and every quotient-functional theta branch containing
\(c\) induce the same quotient map from the source \(J\)-class, by the
argument of \cref{lem:typed-transition-cylinder}.  Hence all direct and
branch-forced transitions stored in the root signature hold in the actual
quotient system.  Their compositions, inverses and upward consequences are
literal identities of that quotient system.  This proves typed root forcing.

Suppose first that \(C\) has a hidden completion.  Compose it with an embedding
of its finite target into a full typed hidden expansion and apply
\cref{lem:typed-real-sort-splitting}.  Use the resulting pairwise distinct
real points as the real witnesses and designated representatives.  Every
profile imposed in \cref{con:typed-guarded-lift} is then satisfied, hence the
actual root pair type belongs to the corresponding exact intersection \(D\).
Insert all private theta interiors by their path-profile realisation and
repeated visible SAP over the already chosen roots.  Support-normality lets us
keep the distinguished complete old higher-support factors fixed.  This gives
a visible completion of \(G_{\mathrm{typ}}(C)\).

Conversely, take a visible completion of \(G_{\mathrm{typ}}(C)\).  Send an
old real vertex to its real witness and a vertex \(b\in B_{a,J}\) to the
actual \(J\)-class of its designated witness.  The first part of the proof
shows that every displayed projection, upward map, quotient transition,
semantic cylinder and old real relation is preserved.

We check reflection, including possible collisions
of unary-term values. On one-generated
functional closures it is exactly the
bounded local diagram check from
\cref{con:typed-guarded-lift}.
For ANY two closures joined in an
irreducible source piece,
\cref{lem:typed-two-root-complete-reflection}
proves that EVERY old atomic
representative pair selected from
their nonempty guarded exact menu
interprets precisely the source's
COMPLETE induced two-generated
hidden type, including the
absence of incompatible named
cylinders and the distinctness
of different term values when
that irreducible join requires
them to be distinct. This
uses the source's full bounded
local reflection check and the
actual semantic-cylinder
partition, not the unjustified
claim that an arbitrary
PARTIAL two-root diagram
determines one unique orbit.
For a support-\(\ge3\)
joining tuple, the old
distinguished complete factor
fixes its canonical pair
shadow, and the separate
support-normal completion
interface must preserve the
higher-support tuples
themselves.

To pass from these bounded ONE-
and TWO-generated checks
to an ARBITRARY irreducible
substructure of the original
partial hidden source, apply
\cref{lem:typed-binary-irreducible-local-cover}.
In the BINARY/unary hidden
signature, any two distinct
vertices of an irreducible
source piece lie together
in some ACTUALLY displayed
one- or two-generated
functionally closed support;
otherwise that piece
admits a nontrivial free
amalgamation decomposition.
Every such bounded support
is embedded inducedly by
the just-verified complete
old local checks. Therefore
the map is injective and
reflects ALL binary named
relations on EVERY
irreducible source
substructure, not merely
on the individual checked
small pieces. It is a
homomorphism-embedding
and hence a hidden
completion of \(C\).
For arity at least three,
the separate support-normal
reflection contract is still
required; pairwise closure
alone is insufficient.
\end{proof}

We next repair factorisation compression.  Let
\[
 S_{\rm typ}=\{0\}\sqcup\bigsqcup_{a,b}\Rch^{ab}
\]
and multiply non-composable profiles to \(0\).  This is a finite semigroup, so
Kufleitner's factorisation-forest theorem applies literally.  A nonzero idempotent belongs to some
\(\Rch^{aa}\) and consists of same-sort
\emph{diversity} atoms. Thus it is irreflexive
as an actual binary relation. The assertion that
the profile \(e\) itself is a reflexive equivalence
would be FALSE. All same-sort atomic colours are
symmetric, so define
\[
                 E_e=\Id_a\cup e.
\tag{TF0}\label{eq:typed-idempotent-equivalence}
\]
This is reflexive and symmetric. For transitivity,
if \(xE_eyE_ez\) and any two of \(x,y,z\)
coincide, the conclusion is immediate; otherwise
the two incident old diversity atoms are in \(e\),
so the STRONG idempotence \(e\starprod e=e\)
forces the diversity colour of \(xz\) into \(e\).
Hence \(E_e\) is an invariant equivalence whose
diversity palette is exactly \(e\).
The following
\cref{lem:typed-idempotent-raw-exact}
strengthens this to a FULL RAW exact
equivalence, making meet-irreducible
reconstruction legal rather than merely
assuming exactness.
Whenever an idempotent profile below is called
an equivalence, the actual relation meant is \(E_e\);
the deleted word still has the strong off-diagonal
profile \(e\), not logical identity.

\begin{lemma}[Literal old-coordinate seams for successive macro sides]
\label{lem:typed-old-seam-telescope}
Let \(C\) be a functionally closed
PARTIAL old hidden structure,
and let \(b_0,\ldots,b_m\)
be a sequence of ACTUAL OLD
hidden source vertex occurrences
with consecutive pairs
\((b_{i-1},b_i)\) issuing
source-supported guarded macro
sides. For each side suppose
\cref{lem:typed-old-source-branch-equation}
provides a literal equation
\[
 \pi_i(\tau_{i-1}(b_{i-1}))
       =\tau_i(b_i)
\]
between NAMED OLD unary-term
values at represented quotient
coordinates \(J_{i-1},J_i\).
Here at a shared intermediate
root \(b_i\) both incident sides
must use the SAME represented
\(J_i\)-coordinate readout,
possibly written by two terms
which are IDENTICALLY the same
canonical quotient map in the
actual hidden signature.

Then their composite equation
\[
 (\pi_m\circ\cdots\circ\pi_1)
    (\tau_0(b_0))=\tau_m(b_m)
\tag{OSE}\label{eq:typed-old-seam-telescope}
\]
holds LITERALLY in \(C\),
regardless of whether the
complete irreducible OLD
supports of two successive
sides have any common REAL
vertex. Their intersection may
consist solely of proper
quotient-coordinate vertices.

The conclusion cannot be
obtained merely because two
DISTINCT OLD hidden source
vertices \(q,q'\) happen to
have EQUAL images in some
homomorphism-embedding completion.
They may be used as one seam
only if their equality is
already literal in \(C\), or
an explicitly displayed NAMED
unary transition in the OLD
source connects their required
coordinate values.
\end{lemma}
\begin{proof}
Because \(C\) is closed under
the named unary functions,
each equation of a macro side
is an equality of actual
vertices of \(C\), not a
constraint on separately
chosen future representatives.
At each intermediate vertex
\(b_i\), both incident term
values denote its same
canonical \(J_i\)-coordinate.
Their equality therefore
holds in \(C\) as well;
if the notations differ by a
quotient transition, the
stated hypothesis demands
that this transition is one
of the actual named unary
maps with its equality already
present in \(C\).
Ordinary substitution and
composition of unary-function
equalities now give (OSE)
by induction on \(m\).
This argument requires no
amalgamation of the distinct
OLD irreducible supports and
no equality of their real
root positions.

For the last assertion,
a homomorphism-embedding of
a reducible partial source
need not be globally
injective. Two distinct old
vertices in different free
amalgamation factors may
acquire the same target image
without being literally
equal source vertices.
Replacing them by one
intermediate term would
therefore add an unproved
OLD equation. The path
contracts through a common
formal OLD root occurrence,
not through an accidental
coincidence of target images.
\end{proof}

\begin{lemma}[A strong idempotent has an exact RAW equivalence envelope]
\label{lem:typed-idempotent-raw-exact}
Let \(a\) be an actual named transitive
sort of a finite complete homogeneous
symmetric binary SAP+semantic-CCO
source. Let \(e\in\Rch^{aa}\) be a
NONEMPTY strong diversity profile
of one positive typed atomic word
\(w:a\to a\), and suppose
\(e\starprod e=e\).
Then
\[
                 E_e=\Id_a\cup e
\]
is not merely an invariant
equivalence: it is an actual
FULL RAW exact equivalence,
represented by the equality-inclusive
endpoint word \(ww\).
Consequently, in the finite
distributive lattice of represented
full exact equivalences on \(a\),
\[
 E_e=\bigcap\{J:J\text{ is a represented
 meet-irreducible equivalence and }
                         J\supseteq E_e\}.
\tag{TF1}\label{eq:typed-idempotent-MI-meet}
\]
The intersection of an empty family
is the universal same-sort equivalence.
Neither a synthetic equality-labelled
diversity edge nor an unproved
exactness of EVERY invariant
equivalence is used.
\end{lemma}
\begin{proof}
Let \(R_w\subseteq M_a^2\) be the
ACTUAL equality-inclusive represented
endpoint relation of the positive word
\(w\). By typed SAP injectivisation
\cref{lem:typed-injectivisation},
its restriction to pairs of DISTINCT
points is exactly the strong
diversity palette \(e\).
Therefore \(R_w\) is either
\(e\) as an actual binary relation
or \(\Id_a\cup e\), depending on
whether the word also has a
genuine equal-endpoint walk.

All invariant a-to-a relations
are unions of symmetric original
same-sort atomic orbitals, hence
\(R_w\) is symmetric. It has an
actual distinct-endpoint witness
because \(e\ne\varnothing\);
by transitivity this gives a
nontrivial \(R_w\)-neighbour for
EVERY point of the sort.
For each real point \(x\), choose
\(y\ne x\) with \(xR_wy\).
Symmetry gives \(yR_wx\), so
\(\Id_a\subseteq R_w;R_w\).

On DISTINCT endpoints, membership
in \(R_w;R_w\) implies membership
in the actual equivalence
\(E_e\) (since \(R_w\subseteq E_e\)
and \(E_e\) is transitive by
the strong idempotence), so the
off-diagonal output lies in \(e\).
Conversely every atom of \(e\)
occurs in \(e\starprod e=e\),
hence has an actual two-step
witness whose edges lie in
\(R_w\); this puts that atom
into \(R_w;R_w\). Thus
\[
                    R_w;R_w
                       =\Id_a\cup e=E_e.
\]
But \(R_w;R_w\) is EXACTLY
the represented RAW endpoint
relation of the concatenated
positive word \(ww\), including
logical identity outputs.
It is a singleton RAW word
profile and therefore a full
exact domain by definition.
It is already an actual
invariant equivalence, so an
element of the represented
full exact-equivalence lattice.

The exact lattice is finite and
distributive by the typed RAW
exact/Helly and lattice criterion
proved earlier in Part~I
(\cref{thm:full-typed-exact-balls,thm:abstract-exact-crt}); every
element of a finite distributive
lattice equals the meet of all
meet-irreducibles above it.
This proves
\eqref{eq:typed-idempotent-MI-meet},
with the ordinary empty-meet
convention when \(E_e=\Top_a\).
\end{proof}

\begin{remark}[Scope of the RAW exact idempotent envelope]
\label{rem:typed-idempotent-envelope-scope}
The preceding equation
\(R_w;R_w=\Id\cup e\)
uses only symmetric original
pair orbitals, source transitivity,
the STRONG endpoint identity
\(e\starprod e=e\), and SAP
injectivisation of RAW walks.
It does NOT require the
terminal no-finite-cover
hypothesis or a priori
exactness of every invariant
equivalence. Its assertion
that \(E_e\) is FULL RAW exact
comes from the literal
positive old word \(ww\).
The subsequent representation
of \(E_e\) as the meet of
represented meet-irreducibles
uses the separate FINITE
DISTRIBUTIVE exact-lattice
theorem from Part~I; this
extra deduction must not
be conflated with the direct
word identity.
\end{remark}

\begin{lemma}[CCO-free idempotent lattice and its distributive CCO specialization]
\label{lem:typed-strong-idempotent-lattice-monolith}
Fix ONE original named TRANSITIVE sort $a$
of a finite-complete homogeneous
binary source whose same-sort old
atom colours are SYMMETRIC.
Assume visible SAP and FULL
original invariant-equivalence NR2
on that sort, but do NOT assume CCO.
Let $\mathscr I_a$ be the set
of all NONEMPTY same-sort
STRONG POSITIVE REACHABLE
word profiles
$e\in\Rch^{aa}$ with
$e\starprod e=e$.
Then
\[
                  e\longmapsto
                    E_e=\Id_a\cup e
\]
is an order-preserving BIJECTION
from $\mathscr I_a$ onto the
finite lattice of ALL original
invariant equivalences
on $a$ EXCLUDING $\Id_a$.
In particular $\mathscr I_a$
is a finite LATTICE with
\[
  e\wedge f=e\cap f\ne\varnothing,
  \qquad
  e\vee f=e\starprod f
             =(E_e\vee E_f)^\circ .
 \tag{IL}\label{eq:typed-strong-idempotent-lattice}
\]
It has a UNIQUE least nonzero
idempotent $e_{0,a}=L_a^\circ$,
where $L_a>\Id_a$ is the
SAP-monolith, and
\[
       e_{0,a}\subseteq e,\qquad
       e_{0,a}\starprod e
                   =e\starprod e_{0,a}=e
       \quad(e\in\mathscr I_a).
\]
If the original source ALSO has
semantic HOMOMORPHIC CCO,
this finite lattice is moreover
DISTRIBUTIVE by original FULL
RAW exact-ball Helly and CRT.
Thus the BIJECTION, NONEMPTY
MEETS and common monolith
hold even in a symmetric
SAP+NR2 source which is
NOT CCO; only distributivity
uses the additional CCO input.

The element $e_{0,a}$ acts as
a unit ONLY inside the
idempotent lattice, NOT on
arbitrary reachable profiles.
This does NOT turn OLD G5
idempotent interval certificates
for the generally COARSER
$E_e$ into an unsupported
$L_a$-equality certificate.
\end{lemma}
\begin{proof}
A nonempty strong idempotent
$e$ gives the ACTUAL original
invariant equivalence
$E_e=\Id_a\cup e$:
symmetry and reflexivity
are immediate, and if
$xE_eyE_ez$ with three
DISTINCT points, the strong
idempotence $e\starprod e=e$
gives $xE_ez$; all coincident
cases are trivial. Thus the
map is injective and takes
values ABOVE $\Id_a$,
without using CCO.

For surjectivity take any
original invariant $E>\Id_a$.
By \cref{lem:typed-all-invariant-equivalences-raw-word}, $E$ is
represented by one old atomic RAW word: FULL NR2 supplies that lemma's
weaker no-relative-index-two premise. Since $E>\Id_a$, the word is
nonempty. This step uses symmetry and homogeneity, but neither SAP nor CCO.
Visible SAP injectivisation
\cref{lem:typed-injectivisation}
turns its true DIVERSITY part
$e=E^\circ$ into a STRONG
positive reachable profile.
Original SAP equivalence
reproduction
\cref{lem:typed-equivalence-reproduction}
gives $e\starprod e=e$.
Thus every $E>\Id_a$ occurs.

For $e,f\in\mathscr I_a$
put $E=E_e,F=E_f$.
Old SAP's invariant monolith
\cref{lem:rh-sap-invariant-monolith}
ensures $E\cap F>\Id_a$.
The NONEMPTY diversity part
$e\cap f=(E\cap F)^\circ$
is itself positive reachable
and idempotent by the preceding
surjectivity, giving the meet.

For the join, apply the two-factor conclusion
\eqref{eq:typed-equivalence-join-strong} of
\cref{lem:typed-raw-strong-saturation} to the actual nonidentity
equivalences $E,F$. Its SAP and same-sort symmetry hypotheses hold here,
and it gives
\[
       E^\circ\starprod F^\circ=(E;F)^\circ=(E\vee F)^\circ.
\]
Thus the join calculation uses no CCO and introduces no STRONG identity atom.

The original SAP-monolith
$L_a>\Id_a$ is the unique
least nonidentity invariant
equivalence by
\cref{lem:rh-sap-invariant-monolith}.
Its positive diversity part
$e_{0,a}=L_a^\circ$ therefore
lies below EVERY strong
idempotent. The product/join
identity gives
$e_{0,a}\starprod e
=(L_a\vee E)^\circ=e$.
Finally, if CCO is ALSO
assumed, terminal full
RAW exactness and the
equivalence-class CRT
\cref{lem:typed-terminal-all-equivalences-exact}
make the FULL invariant
equivalence lattice distributive.
Its nonidentity part is
meet-closed by SAP, so
$\mathscr I_a$ is also
distributive. No synthetic
STRONG equality atom was
ever used.
\end{proof}

\begin{lemma}[Quotient-functional branches are visible at their hidden roots]
\label{lem:typed-branch-root-visibility}
Let \(b,c\) be two distinct hidden root positions,
each either a real root or a proper represented quotient
root. Write \(H_b,H_c\) for the equivalence kernels
of their designated real representatives:
\(H_b=\Id\) for a real root and its original
represented proper meet-irreducible for a quotient
root; similarly for \(c\). Let \(D\) be the
nonempty exact pair menu imposed by the
complete locally valid two-root hidden pattern.
Thus
\[
              H_b;D;H_c=D .
\tag{TV1}\label{eq:typed-anchored-bisaturation}
\]
Suppose a positive typed branch profile \(P\)
in a theta presentation contains \(D\), and
\(P;P^\smile\subseteq J\) for a represented
proper meet-irreducible \(J\) at its source.
Let \(J^P\) be the transported target equivalence.
Then
\[
                  H_b\le J,\qquad
                  H_c\le J^P .
\tag{TV2}\label{eq:typed-branch-visible-kernels}
\]
Consequently the quotient transition forced by
\(P\) is an equality of named unary term values
already in the finite functionally closed
two-root hidden substructure; it does
\emph{not} require naming the \(J\)-coordinate
of an independently chosen, unconstrained
auxiliary representative.
\end{lemma}
\begin{proof}
For any \(xH_bx'\), choose \(y\) with
\(xDy\), possible because \(D\) is nonempty,
invariant, and the source sort is transitive.
By left \(H_b\)-saturation also \(x'Dy\).
Both pairs lie in \(P\), so
\(x(P;P^\smile)x'\), and hence \(xJx'\).
This proves \(H_b\le J\).

The transported equivalence \(J^P\)
contains \(P^\smile;P\), since \(J\)
is reflexive. If \(yH_cy'\), choose
\(xDy\); right \(H_c\)-saturation
gives \(xDy'\). Both pairs lie in \(P\),
so \(y(P^\smile;P)y'\) and hence
\(yJ^Py'\). Therefore \(H_c\le J^P\).
These are inclusions of actual equivalences
and use no same-sort commutativity
of mixed arrows.

For a real root, the named projection
\(F_J\) gives its \(J\)-coordinate.
For a proper quotient root with
\(H_b\le J\), the named upward quotient
gives its source \(J\)-coordinate as
a unary term of \(b\). By
\cref{lem:typed-branch-atomic-transport},
the quotient-functional STRONG branch
\(P\) has the SAME quotient map as
any one of its actual atomic
diversity orbitals \(t\subseteq P\).
Therefore that map, and its target
\(J^P\)-coordinate of \(c\), are
ALREADY named finite-language
unary transition terms; no
new branch-dependent quotient
symbol has to be introduced.
The relation \(D\subseteq P\)
forces their quotient-transition equation
in every actual full hidden expansion.
Since the original induced two-root functional
closure passed the bounded complete-pattern
check of \cref{con:typed-guarded-lift},
the same named term equation already holds
there. Otherwise that two-root closure is
itself a bounded malformed obstruction,
so this branch is never reached.
Thus the transition equation is supported
on the old hidden roots, not on optional
new representative choices.
\end{proof}

\begin{lemma}[Original-source equations for quotient-functional theta branches]
\label{lem:typed-old-source-branch-equation}
Let \((b,c)\) be one ACTUAL
old hidden endpoint pair on which
\cref{con:typed-guarded-lift}
issues a nonempty exact visible
root menu \(D_{bc}\). Let \(H_b,H_c\)
be the original representative kernels.
Suppose \(P\in\Rch^{ab}\) is a
positive branch profile in its
private theta presentation, with
\(D_{bc}\subseteq P\), and let
\(J\) be a represented proper
meet-irreducible with
\(P;P^\smile\subseteq J\).

Then the \(J\)-coordinate of
\(b\) and the transported \(J^P\)
coordinate of \(c\) are NAMED
unary terms of the OLD source
vertices, and the equation
\[
  \pi_{P,J}(\tau_J(b))
           =\tau_{J^P}(c)
\tag{TOE}\label{eq:typed-old-source-branch-equation}
\]
already holds as a LITERAL
unary-term identity in the
original functionally closed
partial hidden structure \(C\).
Here \(\pi_{P,J}\) denotes the
existing named atomic transition
representing the quotient map
of \(P\). The conclusion remains
valid when \(D_{bc}\) is a MEET
of several original profiles,
none of which individually has
\(J\)-functional fibres.

The conclusion applies only to
the source-supported root-pair
gadgets specified in
\cref{con:typed-guarded-lift}.
No such equation is inferred
from an arbitrarily selected
new completion type on an
UNJOINED pair of hidden roots.
\end{lemma}
\begin{proof}
By \cref{lem:typed-private-theta-cycle},
the two gadget endpoints are the
witnesses for actual OLD source
vertices \(b,c\), not artificial
interior roots. By the construction
their positive constraints come
either from one checked
one-generated unary closed
support, or from the complete
induced diagram on a genuinely
irreducible, at most two-generated
OLD source support. The latter
is already a completed hidden
diagram by the bounded check and
\cref{lem:typed-joined-source-support}.
An original higher-support tuple
is treated within its distinguished
complete OLD factor, not by
postulated missing binary data.

For the exact menu \(D_{bc}\),
\cref{lem:typed-branch-root-visibility}
gives
\(H_b\le J\) and
\(H_c\le J^P\).
Thus \(\tau_J(b),\tau_{J^P}(c)\)
are indeed values of NAMED
unary terms in the source
closed support; the quotient
map \(\pi_{P,J}\) is a named
atomic transition by
\cref{lem:typed-branch-atomic-transport}.

Take a realisation of that
COMPLETE old one- or two-root
support in a full hidden source.
Choose distinct representative
points for its old roots, by
\cref{lem:typed-real-sort-splitting}.
By
\cref{lem:typed-one-source-support-per-guarded-pair},
ALL positive spectra on
this witness pair are
pulled back from the SAME
complete irreducible OLD
closed source \(D_{bc}\),
not from separately chosen
local models. Thus its
single actual representative
pair atom \(t\) lies in
EVERY imposed profile
and hence their nonempty
meet \(D_{bc}\).
In particular \(t\in P\).
Since \(P;P^\smile\subseteq J\),
the quotient-transport lemma
forces the displayed equation
\eqref{eq:typed-old-source-branch-equation}
in this realisation.

Crucially, the map from the
old closed source support to
this full hidden realisation
is an EMBEDDING, not merely
a homomorphism; it reflects
equality of the values of
all named unary terms.
The same equation therefore
holds literally on the old
source term values in \(C\).
This last inference uses
COMPLETENESS/IRREDUCIBILITY
of the OLD support, rather
than postulating a term equation
from a new theta interior.
It also uses only the total
meet \(D_{bc}\subseteq P\),
and nowhere requires any
single intersected primitive
profile to be quotient-functional.
\end{proof}

\begin{lemma}[Off-diagonal idempotent interval reinsertion]
\label{lem:typed-idempotent-reinsert}
Let \(w\) be a positive composable typed diversity word
whose STRONG endpoint profile is a nonempty same-sort
idempotent \(e=e\starprod e\). Let \(x,y\) be
\emph{distinct actual real vertices} of its common endpoint
sort. If \(xE_ey\), where \(E_e=\Id\cup e\) is
the actual equivalence of
\eqref{eq:typed-idempotent-equivalence}, then the
entire word \(w\) can be realised between \(x,y\)
with all path interiors new, preserving the
already completed finite old structure on its roots.
\end{lemma}
\begin{proof}
The roots are distinct, hence their actual old
pair colour is a diversity atom. Membership in
\(E_e\) therefore means precisely
\(\col(x,y)\in e\), \emph{not} logical equality.
Since \(e\) is the STRONG endpoint profile of
\(w\), a path realisation with this prescribed
fine endpoint atom exists with its formal path
vertices distinct. Old pair homogeneity places
its two endpoints at the specified pair \(x,y\).
Visible SAP amalgamates the private interior of
that complete finite path with the already
completed finite old structure over their
COMPLETE common two-point induced substructure.
No return to the same real point or synthetic
equality edge is assumed.
\end{proof}

\begin{lemma}[Simultaneous private word insertion over a complete old base]
\label{lem:typed-simultaneous-word-insertion}
Let \(A\) be a finite COMPLETE
visible member of the original
typed SAP+CCO age. Fix a finite
list of old typed diversity
atomic words \(w_i\), each
with two specified DISTINCT
endpoint vertices \(p_i,q_i\in A\)
of the required endpoint sorts.
Suppose
\[
 \col_A(p_i,q_i)\in\Prof(w_i)
\]
for every \(i\), where
\(\Prof(w_i)\) is the genuine
STRONG simple-path endpoint
profile of that particular
original atomic word.
Then a single finite complete
visible member \(B\supseteq A\)
realises ALL words \(w_i\)
between the prescribed endpoints,
with every word's interior
pairwise distinct and with
interiors of different words
disjoint from each other and
from \(A\). Every originally
specified atom on \(A\) is
preserved literally. The
endpoint pairs may overlap
with each other inside \(A\).
\end{lemma}
\begin{proof}
Proceed by finite induction
on the list of obligations.
For \(w_i\), membership
\(\col_A(p_i,q_i)\in\Prof(w_i)\)
means that there is a finite
complete OLD visible structure
\(W_i\) containing one
simple realisation of
the EXACT specified word
\(w_i\) with its endpoints
of the same original fine
atomic pair type as
\((p_i,q_i)\).
By the homogeneous
named-atomic pair type,
identify the two endpoints
of \(W_i\) with \(p_i,q_i\).
Strongly amalgamate the
already constructed complete
base with \(W_i\) over
their complete induced
two-point substructure.
SAP guarantees that all
other vertices of \(W_i\)
remain new, while the full
old base, including every
previously realised word,
is embedded unchanged.
Repeat for all \(i\).
No condition on the NEW
cross colours between
different private interiors
is imposed, so SAP applies
regardless of overlap among
the pinned endpoint pairs.
\end{proof}

